% ──────────────────────────────────────────────────────────────
\subsection{Lorentz covariance}\label{sec:lorentz-covariance-add}

The covariant form~\eqref{eq:Tmunu-dist-def} immediately yields
the transformation law, resolving
Exercise~\ref{ex:Tmunu-tensor}:

\begin{proposition}[Lorentz covariance]\label{prop:lorentz-cov}
  Under a Lorentz transformation
  $x'^\mu = \Lambda^\mu{}_\nu\, x^\nu$, the distributional
  energy-momentum tensor transforms as
  \begin{equation}\label{eq:Tmunu-lorentz}
    T'^{\alpha\beta}(x')
      = \Lambda^\alpha{}_\gamma\,
        \Lambda^\beta{}_\delta\,
        T^{\gamma\delta}(x)\,.
  \end{equation}
\end{proposition}

\begin{proof}
  Under $x' = \Lambda x$, the transformed worldline is
  $\gamma'(\tau) = \Lambda\gamma(\tau)$ and the transformed
  four-momentum is $p'^\alpha(\tau) = \Lambda^\alpha{}_\gamma\,
  p^\gamma(\tau)$.  (Since $\tau$ is Lorentz-invariant and
  $\gamma^\alpha$ transforms as a four-vector,
  $\dot\gamma^\alpha = d\gamma^\alpha/d\tau$ transforms as a
  four-vector, and hence
  $p^\alpha = m\dot\gamma^\alpha$ does as well.)

  We define the transformed distribution~$T'^{\alpha\beta}$ by
  its action on test functions: for any
  $\phi \in \mathcal{D}(\Rn{4})$,
  \begin{align}
    \langle T'^{\alpha\beta},\,\phi\rangle
      &= \sum_j \int_{\R} p_j'^{\,\alpha}(\tau)\,
        \dot\gamma_j'^{\,\beta}(\tau)\,
        \phi\bigl(\gamma_j'(\tau)\bigr)\;\dd\tau
      \notag\\
      &= \sum_j \int_{\R}
        \Lambda^\alpha{}_\gamma\, p_j^\gamma(\tau)\;\cdot\;
        \Lambda^\beta{}_\delta\, \dot\gamma_j^\delta(\tau)\;\cdot\;
        \phi\bigl(\Lambda\gamma_j(\tau)\bigr)\;\dd\tau
      \notag\\
      &= \Lambda^\alpha{}_\gamma\,\Lambda^\beta{}_\delta
        \sum_j \int_{\R} p_j^\gamma(\tau)\,
        \dot\gamma_j^\delta(\tau)\,
        (\phi \circ \Lambda)\bigl(\gamma_j(\tau)\bigr)\;\dd\tau
      \notag\\
      &= \Lambda^\alpha{}_\gamma\,\Lambda^\beta{}_\delta\;
        \bigl\langle T^{\gamma\delta},\;
        \phi \circ \Lambda\bigr\rangle\,.
      \label{eq:Tmunu-lorentz-dist}
  \end{align}
  This is the standard distributional transformation law
  $T'^{\alpha\beta} = \Lambda^\alpha{}_\gamma\,
  \Lambda^\beta{}_\delta\,(\Lambda^{-1})_\ast T^{\gamma\delta}$,
  which is the distributional restatement
  of~\eqref{eq:Tmunu-lorentz}.
\end{proof}

\begin{remark}
  The $3{+}1$ form~\eqref{eq:Tmunu-dt} is \emph{not} manifestly
  covariant: neither $\delta^{(3)}$ nor $d\gamma^\beta/dt$ is
  Lorentz-invariant in isolation.  It is most naturally through
  Theorem~\ref{thm:equivalence}---whose proof rests on the coarea
  formula---that we can deduce that the combination
  $p^\alpha\,(d\gamma^\beta/dt)\,\delta^{(3)}$ transforms as a
  tensor.  The miracle is that the Jacobian factors conspire to
  produce a Lorentz-invariant measure along the worldline; the
  coarea formula explains \emph{why} they conspire.
\end{remark}
